% Generated draft: requires mathematical and source review.
\documentclass[11pt]{article}
\usepackage[margin=25mm]{geometry}
\usepackage{fontspec}
\setmainfont{Latin Modern Roman}
\usepackage{amsmath,amssymb,mathtools,unicode-math}
\setmathfont{Latin Modern Math}
\usepackage{enumitem,needspace,xurl,hyperref}
\setlength{\parindent}{0pt}
\setlength{\parskip}{5pt}
\setlength{\emergencystretch}{3em}
\newcommand{\N}{\mathbb{N}}
\newcommand{\Z}{\mathbb{Z}}
\newcommand{\Q}{\mathbb{Q}}
\newcommand{\R}{\mathbb{R}}
\newcommand{\C}{\mathbb{C}}
\newcommand{\F}{\mathbb{F}}
\newcommand{\A}{\mathbb{A}}
\newcommand{\PP}{\mathbb{P}}
\newcommand{\E}{\mathbb{E}}
\newcommand{\eps}{\varepsilon}
\newcommand{\dd}{\,\mathrm{d}}
\begin{document}
\section*{UnsolvedMath ID: 30001785}
\textbf{Source problem code:} catalogue entry
\subsection{Definitions and mathematical statement}
A standard interval model is \(I=[0,1)\), with its relative Euclidean topology. An interval-exchange transformation is a bijection \(T:I\to I\) for which there is a finite partition of \(I\) into half-open intervals such that \(T\) restricts to a translation on each interval. In this model, \(T\) is minimal if \(\{T^n x:n\geq 0\}\) is dense in \(I\) for every \(x\in I\). A Borel probability measure \(\mu\) is \(T\)-invariant if \(\mu(T^{-1}A)=\mu(A)\) for every Borel set \(A\subseteq I\). Such a measure is ergodic if every Borel set \(A\) satisfying \(\mu(T^{-1}A\mathbin{\triangle}A)=0\) has \(\mu(A)\in\{0,1\}\). The transformation is non-uniquely ergodic if it admits at least two distinct invariant Borel probability measures.

A subset \(R\subseteq I\) is residual if \(I\setminus R\) is meagre, meaning a countable union of nowhere dense sets. One possible precise meaning of saying that \(R\) carries \(\mu\) is
\[
\operatorname{Carry}(\mu,R)\quad\Longleftrightarrow\quad
\text{there exists a Borel set }A\subseteq R\text{ with }\mu(A)=1.
\]
For Borel \(R\), this is equivalent to \(\mu(R)=1\). This convention concerns concentration of a measure; it does not mean that its topological support is contained in \(R\).

The supplied question is: \emph{Can any residual set carry an ergodic measure of a minimal non-uniquely ergodic interval-exchange transformation?} If the residual set is intended to lie in the phase interval, ``any'' means ``every'', the transformation may depend on the set, and ``carry'' has the convention above, the question becomes whether
\[
\begin{gathered}
\forall R\subseteq I\text{ residual},\quad
\exists T:I\to I\ \exists\mu\quad\text{such that}\\
T\text{ is a minimal, non-uniquely ergodic interval-exchange transformation},\\
\mu\text{ is a }T\text{-invariant ergodic Borel probability measure},
\qquad \operatorname{Carry}(\mu,R).
\end{gathered}
\]
This is a conditional formalization, not a resolution of the source's missing qualifications. The supplied record does not specify the ambient space of the residual set, whether the transformation is fixed in advance, the intended quantifier behind ``any'', or the meaning of ``carry''. It also does not explicitly prescribe the interval, endpoint convention, number of exchanged intervals, or measure normalization. Consequently, the displayed interpretation requires human review against the report's surrounding discussion.
\subsection{Short English statement}
The question asks whether residual sets can carry ergodic invariant measures arising from minimal interval-exchange transformations that have more than one invariant probability measure. A possible universal interpretation is displayed, but the ambient space, quantifier order, and meaning of ``carry'' remain unresolved. The supplied record labels the problem open; that label is dated curation, not an independently verified mathematical conclusion.
\subsection{Sources}
\begin{itemize}
\item Supplied local UnsolvedMath registry record, permanent ID \(30001785\), title: \emph{Residual Sets Supporting Interval-Exchange Ergodic Measures}; exact source problem code \texttt{catalogue entry}. Collection: Oberwolfach Reports Open Problems. Proposed year: 2011. No individual proposer is identified in the supplied record.
\item \emph{Billiards, Flat Surfaces, and Dynamics on Moduli Spaces} (2011), Oberwolfach Reports, pp.\ 1361--1427, DOI 10.4171/owr/2011/25. Supplied report reference and literature-source link labeled ``original question'': \url{https://doi.org/10.4171/owr/2011/25}.
\item The supplied background describes Oberwolfach Reports as records of mathematical workshops and their open-problem discussions. The registry marks the statement as exact and reports that its mathematical content was checked against the official report, with notation, context, and question form normalized without changing the source-level claim. This verification description is supplied provenance and has not been independently checked here.
\item Literature assessment dated \textbf{2026-08-21}: the supplied curation labels the problem open, says that the report poses an explicit question, and states that later interval-exchange results provide related residual and ergodic measure theorems without settling this existence question. This is supplied dated curation, not an independently verified literature assessment. The record itself requests independent verification of time-sensitive status and cited claims; no specific later theorem or reference is supplied.
\end{itemize}
\end{document}
